\documentclass[11pt]{article}

\usepackage[utf8]{inputenc}
\usepackage[T1]{fontenc}
\usepackage{lmodern}
\usepackage{geometry}
\usepackage{amsmath,amssymb,amsfonts,amsthm,mathtools}
\usepackage{bm}
\usepackage{enumitem}
\usepackage{microtype}
\usepackage{hyperref}
\usepackage{titlesec}
\usepackage{booktabs}
\usepackage{array}
\usepackage{setspace}
\usepackage{mathrsfs}
\usepackage{physics}

\geometry{margin=1in}
\hypersetup{colorlinks=true, linkcolor=black, urlcolor=blue, citecolor=black}
\setlist[enumerate]{topsep=0.4em,itemsep=0.35em}
\setlist[itemize]{topsep=0.3em,itemsep=0.25em}
\titleformat{\section}{\large\bfseries}{\thesection.}{0.5em}{}
\titleformat{\subsection}{\normalsize\bfseries}{\thesubsection.}{0.5em}{}
\setstretch{1.06}

\newcommand{\Aether}{\AE{}ther}
\newcommand{\Aflow}{\AE{}ther-flow}
\newcommand{\TheoryName}{\Aflow{} Interpretation of Relativity}
\newcommand{\TheoryProgram}{\Aether{} / \Aflow{} framework}
\newcommand{\ObserverMetric}{g^{\mathrm{br}}}
\newcommand{\CandidateMetric}{g^{\mathrm{cand}}}
\newcommand{\CompletedMetric}{g^{\sharp}}
\newcommand{\BaseShift}{\Sigma^{\mathrm{IER}}}
\newcommand{\ResponseShift}{\Sigma^{\mathrm{resp}}}
\newcommand{\CompShift}{\Sigma^{\mathrm{cmp}}}
\newcommand{\ResidualCore}{\mathcal V_{\mu\nu}^{\mathrm{res}}}
\newcommand{\ResidualDec}{\mathcal D_{\mu\nu}^{\mathrm{res}}}
\newcommand{\MixQuad}{\mathcal M_{\mu\nu}^{\mathrm{mix}}}
\newcommand{\RespQuad}{\mathcal Q_{\mu\nu}^{\mathrm{resp}}}
\newcommand{\DeltaTCand}{\Delta T_{\mu\nu}^{\mathrm{cand}}}
\newcommand{\ReducedEin}{\mathcal L_{\ObserverMetric}}
\newcommand{\GaugeBook}{\mathcal G_{\ObserverMetric}}
\newcommand{\CompMix}{\mathcal M_{\mu\nu}^{\mathrm{cmp}}}
\newcommand{\CompQuad}{\mathcal Q_{\mu\nu}^{\mathrm{cmp}}}
\newcommand{\DeltaTComp}{\Delta T_{\mu\nu}^{\mathrm{cmp}}}
\newcommand{\ObsBurden}{\mathfrak B_{\mu\nu}^{\mathrm{obs}}}
\newcommand{\ObserverClass}[1]{\left[#1\right]_{\mathrm{obs}}}
\newcommand{\GaugeExact}{\mathfrak G_{\mu\nu}^{\mathrm{ex}}}
\newcommand{\ConstraintExact}{\mathfrak C_{\mu\nu}^{\mathrm{ex}}}
\newcommand{\RedefTensor}{\mathfrak R_{\mu\nu}^{\mathrm{red}}}

\newtheorem{definition}{Definition}
\newtheorem{proposition}{Proposition}
\newtheorem{remark}{Remark}
\newtheorem{corollary}{Corollary}
\newtheorem{theorem}{Theorem}

\title{\textbf{The \TheoryName{}}\\[0.4em]
\large Substrate-Response / Self-Response Charge-Polarization Candidate\\
\large Exact-Preservation Observer-Equation Exact Absorbability Upgrade or Maximality Note}
\author{}
\date{}

\begin{document}

\maketitle

\begin{abstract}
The immediately preceding one-metric note proved the sharp exact theorem currently available from the fixed exact-preservation package: the displayed observer burden
\[
\ObsBurden
=
\ResidualDec+\CompMix+\CompQuad-8\pi G_N\,\DeltaTComp
\]
is exactly non-independent, so the completed observer theory is exactly observer-side effectively equivalent to Einstein--Hilbert plus ordinary matter through the same operative metric. That theorem still stopped short of the stricter target \(\ObsBurden=0\).

The present note addresses only the remaining upgrade question. It asks whether exact non-independence by itself forces exact absorbability into the same Einstein--matter theory, either through gauge / constraint exact structure or through exact field-redefinition equivalence inside the same one-metric package. The answer is disciplined and two-part. First, the note proves the exact criterion for a successful upgrade: on the fixed package, strict observer-side closure follows if and only if the burden vanishes in the exact observer quotient, equivalently if and only if one has an exact identity placing \(\ObsBurden\) entirely inside already-extracted gauge / constraint exact structure or an exact same-package field-redefinition class that is trivial in that quotient. Second, the note proves that exact non-independence alone does not yield that identity. Therefore the currently recorded package does not yet justify strict observer-side closure, and the preceding effective-equivalence theorem is the strongest honest exact result available on this branch unless genuinely new identities or genuinely new structure are added.

This is a positive maximality theorem for the fixed package, not a universal no-go against every future extension. The note does not claim a completed first-principles substrate derivation of GR.
\end{abstract}

\section{Introduction}

The present manuscript is the narrow follow-up forced by the previous exact non-independence / effective-equivalence note. That note already settled the exact observer-theory status of the displayed burden sectors at the level of independent physics: they do not survive as a second operative metric, a second matter law, a second observer equation, or an unsuppressed linear observer mode. The only remaining observer-side derivational bottleneck on the same one-metric line is therefore sharper and more specific:
\begin{quote}
does exact non-independence upgrade to exact absorbability into the same Einstein--matter theory, or is exact observer-side effective equivalence already the maximal theorem of the fixed package?
\end{quote}

This note answers exactly that question and nothing broader. It does not introduce a deeper origin layer, a modified completion solve, or a second operative metric. It keeps fixed the same candidate metric, the same completed metric, the same completion solve, and the same one-metric matter route, and asks what can honestly be concluded from those inputs alone.

\paragraph{Framework and claim boundary.}
\TheoryProgram{} denotes the broader \Aether{} / \Aflow{} framework built on the claims that the \Aether{} is the underlying four-dimensional substrate of reality, the \Aflow{} is its intrinsic ordered motion, observed three-dimensional space is the local experiential slice of that deeper substrate, S-time is the experienced order of change arising from matter, light, and the \Aflow{}, and observed expansion is the three-dimensional appearance of deeper four-dimensional ordered motion. Within that framework, \TheoryName{} denotes the exact-closure theory in which the effective gravitational dynamics are adopted to be exactly Einsteinian with universal matter coupling. In the active sequence, \emph{adoption} means use of that established relativistic dynamics without claiming substrate derivation, while \emph{derivation} is reserved for a first-principles recovery from explicit substrate variables.

The present note stays strictly inside that boundary. It does not claim that GR has now been derived from substrate microphysics. It does not claim the strict tensor identity \(\ObsBurden=0\). Its positive theorem is narrower: it isolates the exact identity that would be required for that upgrade and proves that, absent such an identity or genuinely new structure, the previous exact observer-side effective-equivalence theorem is already the maximal honest result of the fixed one-metric package.

\section{Fixed Package and Exact Observer Quotient}

\begin{definition}[Fixed one-metric exact-preservation package]
\label{def:fixed}
The present note keeps the following observer-level data fixed.
\begin{enumerate}
    \item The candidate and completed operative metrics are
    \begin{equation}
    \CandidateMetric
    =
    \ObserverMetric+\BaseShift+\ResponseShift,
    \qquad
    \CompletedMetric
    =
    \ObserverMetric+\BaseShift+\ResponseShift+\CompShift.
    \label{eq:metrics}
    \end{equation}
    \item The fixed completion solve is
    \begin{equation}
    \ReducedEin(\CompShift)=-\ResidualCore.
    \label{eq:solve}
    \end{equation}
    \item The completed observer equation is
    \begin{equation}
    G_{\mu\nu}[\CompletedMetric]
    =
    8\pi G_N\,T_{\mu\nu}^{\mathrm{matter}}[\CompletedMetric,\psi]
    +\GaugeBook(\CompShift)_{\mu\nu}
    +\ObsBurden,
    \label{eq:completed}
    \end{equation}
    with
    \begin{equation}
    \ObsBurden
    \equiv
    \ResidualDec+\CompMix+\CompQuad-8\pi G_N\,\DeltaTComp.
    \label{eq:burden}
    \end{equation}
    \item The ordinary matter route is still exactly the same one-metric route
    \begin{equation}
    \DeltaTComp
    \equiv
    T_{\mu\nu}^{\mathrm{matter}}[\CompletedMetric,\psi]
    -
    T_{\mu\nu}^{\mathrm{matter}}[\CandidateMetric,\psi].
    \label{eq:deltatcomp}
    \end{equation}
    \item The inherited candidate-branch decomposition
    \begin{equation}
    \ResidualDec
    =
    \MixQuad+\RespQuad-8\pi G_N\,\DeltaTCand
    \label{eq:resdec}
    \end{equation}
    is kept fixed.
    \item Gauge bookkeeping is already extracted explicitly in \eqref{eq:completed}, and the constraint-exact elimination sectors have already been removed before \(\ObsBurden\) is defined.
    \item The previously established exact theorem on the same branch is that \(\ObsBurden\) is exactly non-independent and the completed branch is exactly observer-side effectively equivalent to Einstein--Hilbert plus ordinary matter.
\end{enumerate}
\end{definition}

\begin{definition}[Exact observer quotient]
\label{def:quot}
For the fixed package of Definition \ref{def:fixed}, the \emph{exact observer quotient} is the quotient of observer-side tensors by the already-recognized null classes:
\begin{enumerate}
    \item gauge-exact bookkeeping tensors;
    \item constraint-exact elimination tensors;
    \item exact same-package field-redefinition tensors whose effect is only to rewrite the same one-metric Einstein--matter equation without changing the operative metric or matter law.
\end{enumerate}
The observer-quotient class of a tensor \(X_{\mu\nu}\) is denoted \(\ObserverClass{X_{\mu\nu}}\).
\end{definition}

\begin{remark}
Definition \ref{def:quot} records the only exact null classes relevant to the present upgrade question. The current note does not enlarge them. In particular, exact non-independence by itself is \emph{not} added as a quotient relation; otherwise the previous theorem would trivialize the present problem by definition rather than by proof.
\end{remark}

\begin{definition}[Exact absorbability on the fixed package]
\label{def:absorb}
The burden \(\ObsBurden\) is \emph{exactly absorbable} on the fixed package if its observer-quotient class vanishes:
\begin{equation}
\ObserverClass{\ObsBurden}=0.
\label{eq:quotvanish}
\end{equation}
Equivalently, \(\ObsBurden\) is exactly absorbable if there exist tensors \(\GaugeExact\), \(\ConstraintExact\), and \(\RedefTensor\) belonging respectively to the three null classes of Definition \ref{def:quot} such that
\begin{equation}
\ObsBurden=\GaugeExact+\ConstraintExact+\RedefTensor.
\label{eq:absorbdecomp}
\end{equation}
\end{definition}

\begin{remark}
Definition \ref{def:absorb} is the precise upgrade target. It is weaker than demanding the literal tensor identity \(\ObsBurden=0\), but it is stronger than exact non-independence. It asks whether the burden is trivial in the only exact observer quotient already licensed by the fixed one-metric package.
\end{remark}

\section{What an Exact Upgrade Would Require}

\begin{proposition}[Exact absorbability is equivalent to exact observer-quotient vanishing]
\label{prop:equiv}
On the fixed package, Definition \ref{def:absorb} is equivalent to either of the following statements:
\begin{enumerate}
    \item \(\ObserverClass{\ObsBurden}=0\);
    \item there exists an exact identity of the form \eqref{eq:absorbdecomp}.
\end{enumerate}
\end{proposition}

\begin{proof}
This is immediate from the definition of a quotient: the class \(\ObserverClass{\ObsBurden}\) vanishes if and only if \(\ObsBurden\) differs from zero by an element of the subspace generated by the declared null classes. Definition \ref{def:quot} names those null classes explicitly, and Definition \ref{def:absorb} records their summed representative as \eqref{eq:absorbdecomp}.
\end{proof}

\begin{proposition}[Exact absorbability upgrades the completed observer equation]
\label{prop:upgrade}
If \(\ObsBurden\) is exactly absorbable on the fixed package, then the completed observer equation is exactly equivalent in the observer quotient to the same Einstein--matter equation carried by \(\CompletedMetric\):
\begin{equation}
\ObserverClass{G_{\mu\nu}[\CompletedMetric]-8\pi G_N\,T_{\mu\nu}^{\mathrm{matter}}[\CompletedMetric,\psi]}=0.
\label{eq:quotein}
\end{equation}
If, in addition, the field-redefinition class in Definition \ref{def:quot} is trivial on the fixed package, then exact absorbability reduces to the stronger criterion
\begin{equation}
\ObsBurden=\GaugeExact+\ConstraintExact.
\label{eq:noredef}
\end{equation}
\end{proposition}

\begin{proof}
Substitute \eqref{eq:absorbdecomp} into \eqref{eq:completed}. The right-hand side beyond the Einstein--matter term is then a sum of tensors whose observer-quotient class is zero by Definition \ref{def:quot}. Hence \eqref{eq:quotein} follows. If the same-package field-redefinition null class is trivial, then \(\RedefTensor=0\) in \eqref{eq:absorbdecomp}, leaving exactly \eqref{eq:noredef}.
\end{proof}

\begin{theorem}[Exact same-package closure criterion]
\label{thm:criterion}
On the fixed one-metric exact-preservation package, strict observer-side closure on the same Einstein--matter branch can be upgraded only by an exact identity proving that \(\ObsBurden\) is trivial in the exact observer quotient. Equivalently, one needs one of the following:
\begin{enumerate}
    \item an exact gauge / constraint absorption identity for \(\ObsBurden\);
    \item an exact same-package field-redefinition theorem whose remainder class is trivial in the observer quotient;
    \item the special stronger identity \(\ObsBurden=0\).
\end{enumerate}
Without one of these exact identities, the completed observer equation does not upgrade to strict same-package closure.
\end{theorem}

\begin{proof}
Propositions \ref{prop:equiv} and \ref{prop:upgrade} show that exact absorbability is exactly equivalent to the vanishing of \(\ObserverClass{\ObsBurden}\), and that this vanishing is what upgrades the completed equation to exact Einstein--matter form in the observer quotient. The special case \(\ObsBurden=0\) obviously implies quotient vanishing. Conversely, if none of the three listed exact identities is available, then \(\ObserverClass{\ObsBurden}\) has not been shown to vanish, so the upgrade has not been proved.
\end{proof}

\section{Why Exact Non-Independence Alone Does Not Force the Upgrade}

\begin{definition}[Exact non-independence, recalled]
\label{def:nonind}
On the fixed package, a tensor sector is \emph{exactly non-independent} if it is an exact functional only of already fixed branch data, introduces no second operative metric, no second matter law, no extra observer equation, and no independent unsuppressed linear observer mode.
\end{definition}

\begin{remark}
Definition \ref{def:nonind} is recalled only to separate the previous theorem from the present stronger upgrade target. Exact non-independence is a statement about independent observer-side content. Exact absorbability is a statement about triviality in the exact observer quotient.
\end{remark}

\begin{proposition}[Exact non-independence does not imply observer-quotient triviality]
\label{prop:notenough}
On the fixed package, the statement that \(\ObsBurden\) is exactly non-independent does not by itself imply \(\ObserverClass{\ObsBurden}=0\).
\end{proposition}

\begin{proof}
The hypotheses are logically different. Exact non-independence constrains how \(\ObsBurden\) depends on branch data and what new observer-side content it fails to introduce. It does not state that \(\ObsBurden\) is gauge exact, constraint exact, or an exact same-package field-redefinition tensor. Therefore exact non-independence alone does not place \(\ObsBurden\) in any of the null classes of Definition \ref{def:quot}. Since quotient triviality requires exactly that membership, the implication does not follow.
\end{proof}

\begin{proposition}[Re-expansion origin is not yet an exact cancellation identity]
\label{prop:reexp}
Treating \(\ResidualDec\), \(\CompMix\), \(\CompQuad\), and \(\DeltaTComp\) as exact re-expansion sectors of the same Einstein--matter package does not by itself prove
\begin{equation}
\ResidualDec+\CompMix+\CompQuad-8\pi G_N\,\DeltaTComp=0
\label{eq:reexpzero}
\end{equation}
or any weaker identity implying \(\ObserverClass{\ObsBurden}=0\).
\end{proposition}

\begin{proof}
The fixed package records only that the displayed sectors are exact functionals of the same branch data and the same one-metric matter route. That information explains why the sectors are not new independent observer-side laws, but it does not yet supply a tensor identity relating their sum to an already-extracted gauge / constraint exact tensor or to a trivial same-package field-redefinition class. In particular, the decomposition \eqref{eq:burden} is an exact bookkeeping identity for the completed equation, not already an exact cancellation identity. Therefore the mere fact that the sectors arise from one common Einstein--matter package does not prove \eqref{eq:reexpzero} or quotient triviality.
\end{proof}

\begin{theorem}[Fixed-package maximality theorem]
\label{thm:max}
Within the fixed one-metric exact-preservation package of Definition \ref{def:fixed}, exact observer-side effective equivalence is the strongest honest exact theorem currently justified. A stronger strict same-package closure theorem would require genuinely new input, namely at least one of:
\begin{enumerate}
    \item a new exact cancellation identity among \(\ResidualDec\), \(\CompMix\), \(\CompQuad\), and \(\DeltaTComp\);
    \item a new proof that their sum is already gauge exact or constraint exact in the completed observer quotient;
    \item a new exact same-package field-redefinition theorem making the burden quotient-trivial.
\end{enumerate}
\end{theorem}

\begin{proof}
The previous note establishes exact non-independence and observer-side effective equivalence on the fixed package. Proposition \ref{prop:notenough} shows that exact non-independence alone does not imply quotient triviality. Proposition \ref{prop:reexp} shows that common re-expansion origin alone does not supply the missing cancellation identity. Theorem \ref{thm:criterion} identifies quotient triviality as the necessary and sufficient upgrade target. Therefore, absent one of the new inputs listed above, the fixed package does not justify a theorem stronger than the existing effective-equivalence result.
\end{proof}

\begin{corollary}[No strict observer-side closure theorem from the recorded package alone]
\label{cor:noclosure}
At current manuscript scope, the recorded fixed package does not yet prove the strict observer-side closure theorem
\[
G_{\mu\nu}[\CompletedMetric]
=
8\pi G_N\,T_{\mu\nu}^{\mathrm{matter}}[\CompletedMetric,\psi]
\]
even modulo the exact observer quotient.
\end{corollary}

\begin{proof}
By Theorem \ref{thm:criterion}, such a theorem would require the vanishing of \(\ObserverClass{\ObsBurden}\). By Theorem \ref{thm:max}, that vanishing has not been justified from the recorded package alone.
\end{proof}

\section{Project-Level Verdict on the Next Move}

\begin{theorem}[Exact upgrade-or-maximality verdict]
\label{thm:verdict}
For the one-metric charge-polarization branch at current scope, the next honest primary move is exactly the following alternative:
\begin{enumerate}
    \item either prove an exact identity that makes \(\ObsBurden\) trivial in the exact observer quotient, thereby upgrading observer-side effective equivalence to strict same-package closure;
    \item or record the maximality verdict of Theorem \ref{thm:max}, namely that the fixed package already yields the strongest honest exact observer-side theorem available without new identities or new structure.
\end{enumerate}
Another deeper-origin manuscript that leaves the exact observer equation unchanged does not answer this remaining upgrade question.
\end{theorem}

\begin{proof}
The first item is exactly Theorem \ref{thm:criterion}. The second item is exactly Theorem \ref{thm:max}. Since the present obstruction is now the observer-quotient status of \(\ObsBurden\), a further deeper-origin note that preserves the same completed observer equation and the same fixed package would be observer-equation neutral with respect to this bottleneck.
\end{proof}

\begin{corollary}[Claim boundary]
\label{cor:boundary}
The present note does not prove:
\begin{enumerate}
    \item the strict tensor identity \(\ObsBurden=0\);
    \item a completed exact field-redefinition equivalence theorem eliminating \(\ObsBurden\);
    \item a completed first-principles substrate derivation of GR.
\end{enumerate}
Its positive claim is narrower and exact: the previous effective-equivalence theorem is maximal for the recorded fixed package unless new identities or new structure are added.
\end{corollary}

\begin{proof}
Items (1) and (2) are excluded by Theorems \ref{thm:criterion} and \ref{thm:max}. Item (3) is excluded by the project-wide adoption / derivation boundary stated in the introduction.
\end{proof}

\section{Conclusion}

The remaining derivational bottleneck on the same one-metric branch is now extremely narrow. The previous note already proved that the displayed burden sectors do not survive as additional independent observer-side physics. The question left open was whether that exact non-independence automatically upgrades to exact absorbability inside the same Einstein--matter package.

The present note shows that the answer is not automatic. The successful upgrade criterion is exact observer-quotient triviality of \(\ObsBurden\): one must prove an exact gauge / constraint absorption identity, an exact same-package field-redefinition theorem with trivial quotient class, or the special stronger identity \(\ObsBurden=0\). Exact non-independence alone does not provide that result, and the fact that the displayed sectors arise as re-expansion terms of the same Einstein--matter package does not by itself yield the missing cancellation identity.

Therefore the current one-metric exact-preservation package has reached a clean maximality point. Its strongest honest exact observer-side result is the already-recorded exact effective-equivalence theorem, not strict same-package closure. The next honest burden is correspondingly sharp: either find the exact identity that kills \(\ObserverClass{\ObsBurden}\), or stop pressing this branch for more than it currently contains and treat the effective-equivalence theorem as maximal unless genuinely new identities or genuinely new structure are added.

\end{document}
